
avec @@M0@@ un @@M1@@-faisceau et @@M2@@. Alors
@@M3@@
mais
@@M4@@
par récurrence sur @@M5@@ et @@M6@@; d'où (i) par récurrence sur @@M7@@. Enfin, par 1.7 et récurrence sur @@M8@@, @@M9@@ et @@M10@@ est un @@M11@@-faisceau; d'où (ii).

\textbf{Définition 1.9.} \emph{On appelle \emph{@@M12@@-polynômes} les polynômes à coefficients rationnels @@M13@@ suivants, définis par récurrence}~:
@@M14@@

\textbf{Remarque 1.10.} Pour les @@M15@@-polynômes, on voit tout de suite par récurrence sur @@M16@@ que
@@M17@@

\textbf{Théorème 1.11.} \emph{(Le théorème principal sur les @@M18@@-faisceaux). Soient @@M19@@ un @@M20@@-faisceau sur @@M21@@ et @@M22@@ le polynôme de Hilbert de @@M23@@; soient @@M24@@ une suite d'entiers @@M25@@ et @@M26@@ l'entier donné par le @@M27@@-ième @@M28@@-polynôme (1.9). Alors @@M29@@ et @@M30@@ est @@M31@@-régulier. En particulier, si @@M32@@, alors @@M33@@ est @@M34@@-régulier.}

